% !TEX root = ../diary-book.tex
\section{令和7年10月01日 全体会議での発表準備（シャワー粒子の頻度について）}
\label{sec:diary-20251001}

この一週間ほどシミュレーションや測位についての検討を重ねてきた。
位置をある程度振ったシミュレーションも必要であるため、多少時間はかかってもやるべきとの考えに至った。

ある場所でミューオンを測定した場合に、そのミューオンを地上の検出器で計測できる確率および立体角で決まっている。
地表面を 100分の位置の面積で覆っている場合には確率は 0.1 となる。
一方でシャワー面の測定はより複雑であり、検出器で３個以上のシャワー粒子を同時計測した場合に同定できたと仮定する。
立体角の割合の３乗で減っていくが一方でシャワー形成が進んでいる場合には粒子の数N の 3乗程度で確率が増加するため、
エネルギーが高くなってくるとシャワー形成が著しじるしく進むことがわかる。

１個の粒子の測定確率が $\epsilon$ のときに、N 個の粒子のうち３個測定できる確率は以下の様に表すことができる。

\begin{equation}
  \epsilon_{\mathrm eff} = _{n}C_{3} \epsilon^3 ~ n^3/6 \epsilon^3
\end{equation}

一方で、N 個の粒子のうち一つを検出する確率は以下の様に表される。
\begin{equation}
  \epsilon_{\mathrm single} = n \epsilon
\end{equation}

ここで注意しなければならないのが、シャワー計測の場合はガンマ線も用いること、直接測定ではミューオンだけを使っていることである。

% !TEX root = ../diary-book.tex

\section{令和7年10月02日 全体会議での発表}
\label{sec:diary-20251002}

話をしたが全く盛り上がることなくなんだかよくわからない発表になってしまった。それでも期待されているとのことでもっと取り組む時間を増やしていく必要がある。
特にどのようなアルゴリズムで処理をしようとしているかをデモンストレートする必要がある。

高橋忠幸さんからは、どのような論文にしようとしているのかを教えてほしいと宿題を与えられた。
今日の発表で他に出た質問は、
\begin{itemize}
  \item どのようなアプリケーションがあるか（潜水艦では磁場マップを用いた量子計測、電波の届く場合は GPS 以外の大体電波、光格子など）とういことについて、その他の手段があることも前提としながら意識を向けるべき
  \item やっぱりどういう仕組で測位・時刻同期しようとしているかわからない
  \item どのくらいの測位精度でやろうとしているのかわからない
  \item 
\end{itemize}

シャワー面の理解をすることはだんだんできてきたので、定量化してまとめを作成していく

\begin{itemize}
  \item 一次粒子の位置および角度分布をいれる（一次宇宙線のフラックス単位は $ ({\mathrm GeV}\cdot m^{2} \cdot s \cdot sr)^{-1} $ となっている）。一次粒子のエネルギーについては累積度数分布に影響を与える可能性のある 10 GeV 以上を対象とすればよい。位置については平均自乗半径の３倍程度まで（4 km 程度）を取っておけばよいだろう。角度については角度分布がどの様になっているかを知っておく必要がある。一次粒子の情報を保存するための仕組みが必要である。一次粒子の位置についてはあとで手で入れればよいので角度だけをシミュレーションに入れておけば良い。エネルギーは相変わらず離散的に扱うことにする。
  \item 一次粒子としてヘリウムを追加する。
  \item ミューオンの到来頻度（手のひら一秒を再現する）。一次粒子の角度分布を入れたうえで、位置を振りながら２次元平面上に
  \item シャワー面の可視化を行い、エネルギーごとの曲率や面の厚さを導出する
  \item シャワー面の同定を真面目にやってみる
\end{itemize}


論文は低エネルギー一次宇宙線に対するシャワー面同定手法の確立と、シャワー面同定を用いた位置導出手法についてに分割し、前者はシャワー観測分野における科学論文として、後者は測位・時刻同期についての論文として出版することができるとよい。

% !TEX root = ../diary-book.tex

\section{令和7年10月6日 角度および地中への入射も含めたシミュレーション}
\label{sec:diary-20251006}

\subsection{一次宇宙線の角度分布について}

まずは一次宇宙線の角度分布について調べる必要はあり、地場を考えず等方的的であるとして、取り扱うことにする。状況を簡単にするために極角方向にはサイコロは振らずに、方位角方向にのみサイコロを降ることにする。$\cos \theta$ で降ることと、立体角が 1 sr 程度になるようにしたい。
COSMOS X では方位角を $\cos \theta$ で示しているが、どこから見て。

